% Ad Familiares 7.12 --- headnote
% ad-familiares-07-12, early February 53 BC, Rome

\textit{Cicero to C. Trebatius Testa in Gaul, written from Rome in
the early days of February 53 BC. Vibius Pansa, returning from
Caesar's camp, has reported to Cicero that Trebatius has gone over
to Epicureanism; Cicero, mock-scandalised, writes at once. The
joke is a triple one: at the camp at Samarobriva (Caesar's winter
quarters in Belgic Gaul) for breeding philosophers fit only for
seaside Tarentum; at Trebatius, who has apparently picked up his
new doctrine from Pansa's circle; and at the friend ``Zeius,''
whose name in the manuscripts is corrupt but who was evidently
another Epicurean of Cicero's acquaintance, sufficient to make
Trebatius's earlier orthodoxy already a little suspect.}

\textit{The second paragraph is sustained legal-philosophical
banter. Epicureanism, with its principle that pleasure is the
measure of all things and that the wise man should not enter
public life, is set against the formulae of the civil law in
which Trebatius is a recognised expert: the action on a
fiduciary trust, with its standard phrase \textit{ut inter bonos
bene agier oportet}; the action \textit{communi dividundo} for
partitioning common property; and the archaic oath \textit{per
Iovem lapidem}, sworn over a stone with imprecations on perjury
--- all of them, Cicero says, drained of force on Epicurean
premises. The final swipe is at the wretched Latin town of
Ulubrae, whose population (later the buzzing frogs of 7.18)
will have no civic life left to them if Trebatius adopts the
Epicurean \textit{lathe biosas}. The Greek verb
\textit{politeuesthai} --- ``to take part in public life'' ---
clinches the point. Cicero closes with a half-serious request
for news.}
